\subsection{The mixed ordinary Hilbert seed at the original level}
\label{sec:mixed-hilbert-seed}

Here is a direct construction of the mixed rank-two seed used in the
second rank-four case. The key is that, on a totally definite
quaternion algebra, ordinary reduction depends only on the residual
weight character. No assertion about the rank-four deformation ring
follows merely from the existence of this seed.

\begin{theorem}[Mixed ordinary weight change]\label{thm:mixed-weight-change}
Let \(F\) be totally real of even degree and let the odd prime \(p\)
split completely in \(F\). Let \(\xi\) be a totally odd finite-order
Galois character, unramified at \(p\), of conductor \(\mathfrak n\).
Suppose that \(\mathfrak n\) is divisible by a prime of norm \(2\).
Use arithmetic Frobenius and write \(\epsilon\) for the \(p\)-adic
cyclotomic character. Fix a nonempty proper subset \(P\) of the primes
above \(p\), and suppose
\[
 \bar\xi^2\ne1,\qquad
 \bar\xi(\operatorname{Frob}_v)\ne1\quad(v\mid p).
\]
Suppose there is a parallel weight \(k+1\) cuspidal Hilbert eigenform
of character \(\xi\), level \(\mathfrak n\prod_{v\mid p}v\), where
\(k\ge2\) is divisible by \(p-1\), with the following residual
eigenvalues in the usual arithmetic Hecke normalization:
\[
 \overline{T_{\mathfrak q}}=1+\bar\xi(\mathfrak q)
 \quad(\mathfrak q\nmid p\mathfrak n),\qquad
 \overline{U_v}=
 \begin{cases}1&v\in P,\\ \bar\xi(v)&v\notin P.\end{cases}
\]
Then for every \(t\ge2\) divisible by \(p-1\) there is a continuous
representation \(r:G_F\to\mathrm{GL}_2(\overline{\mathbb Q}_p)\),
irreducible on every open subgroup, such that
\[
 \det r=\epsilon^k\xi,\qquad
 \bar r^{\rm ss}=1\oplus\bar\xi,\qquad
 r|_{I_{\mathfrak q}}\simeq1\oplus\xi|_{I_{\mathfrak q}}
       \quad(\mathfrak q\nmid p).
\]
At \(v\mid p\), \(r\) is crystalline and has a stable line and
quotient with cyclotomic inertia exponents, respectively,
\[
 (k+t,-t)\quad(v\in P),\qquad (k,0)\quad(v\notin P).
\]
The characters on these lines and quotients reduce respectively to
\((\bar\xi,1)\) and \((1,\bar\xi)\) on the entire decomposition
group. The automorphic attachment can be taken spherical at \(p\)
and outside \(\mathfrak n\), and principal series with distinct
inertial characters at primes dividing \(\mathfrak n\).
\end{theorem}

\begin{proof}
Let \(B/F\) be the quaternion algebra ramified exactly at the real
places; it exists because the degree is even. Fix its splittings at
finite places. Jacquet--Langlands transfers the initial cusp form
without changing its finite local representations. Put
\(K^p=K_0(\mathfrak n)^p\) and at \(v\mid p\) take the upper
Iwahori \(I_v\subset\mathrm{GL}_2(\mathbb Z_p)\). Set
\(K=K^p\prod_{v\mid p}I_v\). At the tame level use the character
of the lower-right entry prescribed by the initial form (the usual
nebentypus type). Write \(O\) for a sufficiently large coefficient
DVR, \(\lambda\) for its uniformizer, and \(\kappa=O/\lambda\).

We first give the integral weight calculation. Put \(m=k-1\).
For integers \(n_v\ge1,b_v\) with \(n_v+2b_v=m\), let
\[
 L=\bigotimes_{v\mid p}
       \operatorname{Sym}^{n_v}O^2\otimes\det^{b_v},
 \qquad \tau=\bigotimes_{v\mid p}\tau_v .
\]
The common central exponent is \(m\). The ordinary algebraic
modular forms on \(B\) have a lattice \(M_L\) of functions on
\(B^\times(\mathbb A_{F,f})\) with left \(B^\times(F)\)-invariance
and right transformation
\[
 f(gzu)=\omega(z)\nu(u^p)\tau(u_p)^{-1}f(g).
\]
Here \(\nu\) is the fixed tame type and \(\omega\) is the
\(p\)-adic avatar of the original automorphic central character:
if the latter is \(\omega_0\), then
\(\omega(z)=\omega_0(z)\prod_{v\mid p}z_v^{-m}\).
In the arithmetic Hecke normalization \(\omega_0=\xi_f|\cdot|_f^{-m}\),
where \(\xi_f\) denotes the finite idelic character giving the
arithmetic Hecke values of \(\xi\). The avatar has unit values:
the contribution of a uniformizer at \(p\) to
\(|\cdot|_f^{-m}\) is canceled by the algebraic factor, and at other
places it is a \(p\)-adic unit. It is trivial on \(F^\times\):
the product of the norm and algebraic factors there is the sign
of the rational norm to the power \(m\), which cancels the totally
odd character. On the intersection of the center and \(K\) the
two right actions agree. Thus the same central character and tame
type are used at every weight. The description over the coefficient
field is the usual algebraic modular form description, obtained by
moving the algebraic action at infinity to the chosen places above
\(p\). No change of level or of central character is involved.

The projective double coset set
\[
 B^\times(F)\backslash B^\times(\mathbb A_{F,f})/
       (Z(\mathbb A_{F,f})K)
\]
is finite. For each representative its rational stabilizer modulo
\(F^\times\) is finite, by discreteness of the projective rational
group in the adeles and compactness at infinity and at the indicated
finite level. At the prime \(u\mid\mathfrak n\) of norm \(2\) this
stabilizer injects, after conjugation by the local representative,
into the projective image of an Iwahori in
\(\mathrm{GL}_2(F_u)\). This is a pro-\(2\) group: the upper
Borel of \(\mathrm{GL}_2(\mathbb F_2)\) has order \(2\), and the
principal congruence kernel is pro-\(2\).
Consequently the finite stabilizer has order a power of \(2\).
Since \(p\) is odd, taking its invariants is exact over \(O\)
and \(\kappa\), by the averaging idempotent. Evaluation at the
coset representatives expresses \(M_L\) as the direct sum of these
invariant modules; in particular reduction commutes with taking
the space, and a surjection of residual coefficient modules gives
a surjection of the residual spaces. This argument includes
\(p=3\) and does not require an auxiliary neat tame level.

In the \(v\)-factor of \(L/\lambda L\), extract the coefficient of
\(e_2^{n_v}\), including the determinant character. The resulting
one-dimensional quotient has Iwahori character
\[
 \begin{pmatrix}a&c\\0&d\end{pmatrix}
   \longmapsto d^{n_v}(ad)^{b_v}.
\]
Let \(q_L\) be the tensor product of these quotient maps. If the
exponents \(n_v,b_v\) are fixed modulo \(p-1\), the scalar quotient
\(L_{\rm sc}\) is independent of the weights. The pointwise map
\[
 \overline M_L\longrightarrow M_{L_{\rm sc}}
\]
is surjective by the stabilizer calculation.

The right cosets of \(I_v\operatorname{diag}(p,1)I_v\) can be
represented by \(h_c=\left(\begin{smallmatrix}p&c\\0&1\end{smallmatrix}\right)\),
\(c\bmod p\). With the above right transformation rule the
normalized operator is
\[
 U_v^{\rm norm}f(g)=p^{-b_v}\sum_{c\bmod p}\tau_v(h_c) f(gh_c).
\]
It preserves \(M_L\), since the factor from \(\det(h_c)^{b_v}\)
cancels the normalization. On the residual symmetric power the
individual matrix sends a vector \(x\) to
\[
 (c e_1+e_2)^{n_v}\,[e_2^{n_v}]x .
\]
Its scalar quotient is just \([e_2^{n_v}]x\). Thus the map to
the scalar space commutes with the normalized \(U_v\), acting
there by the sum over the same \(h_c\), and with Hecke operators
away from \(p\). The scalar operator is well-defined by
surjectivity and this identity. The product \(U=\prod_{v\mid p}U_v^{\rm norm}\)
kills the entire kernel of the scalar quotient modulo \(\lambda\):
each term in its expansion factors through \(q_L\), evaluated
at the corresponding translated argument. Hence the invertible
parts for \(U\) are isomorphic for all these weights. For clarity,
over the residue field the invertible part is the stable image
of a sufficiently high power of \(U\); restricting a surjection
to stable images is surjective, and a kernel killed by \(U\)
has zero intersection with the invertible part. Over \(O\) the
corresponding direct summand is obtained by the ordinary idempotent
(the limit of \(U^{j!}\)); it reduces to that stable image.
This follows also by lifting the coprime factors of the
characteristic polynomial of \(U\) that reduce to powers of \(X\)
and polynomials prime to \(X\). Since the \(U_v^{\rm norm}\) commute,
each is invertible on the ordinary part.

Take \((n_v,b_v)=(k-1,0)\) first, and then take
\[
 (n_v,b_v)=
 \begin{cases}
 (k+2t-1,-t)&v\in P,\\
 (k-1,0)&v\notin P .
 \end{cases}
\]
Their scalar quotients are identical. The initial eigenpacket
has nonzero ordinary reduction, and so it occurs in the ordinary
part of the new integral space modulo \(\lambda\). The
Deligne--Serre lifting lemma for the commuting integral Hecke
operators yields a characteristic-zero eigenpacket with the same
residual values. One can apply it to the finite \(O\)-algebra
generated by the operators in the endomorphism ring of the
ordinary lattice, localize at the given residual maximal ideal,
and choose a characteristic-zero prime below it; the algebra
acts faithfully on a free lattice and has no \(O\)-torsion.
Jacquet--Langlands transfers the resulting representation back
to a cuspidal \(\mathrm{GL}_2\) representation \(\pi_0\):
one-dimensional quaternionic representations cannot occur at
the nonzero symmetric-power weights \(n_v\).
The finite central character and tame level are unchanged.

We check the local conditions. At \(v\mid p\) the representation
is Iwahori-spherical with unramified central character. Its local
generic constituent is either an unramified principal series or
an unramified twist of Steinberg. In the latter case the raw
Iwahori Hecke eigenvalue squared is the central value at a
uniformizer, up to a unit; its valuation is \(m/2\).
The normalized slope would be \(m/2-b_v=n_v/2>0\), contrary
to ordinarity. Hence it is spherical. If \(\alpha_v\) is the raw
ordinary Hecke root, its valuation is \(b_v\) and that of the
other root is \(n_v+b_v+1\), since the determinant of the
arithmetic Hecke polynomial is \(p^{m+1}\xi(v)\).

Use the classical Galois representation and local-global
compatibility theorem for cohomological cuspidal Hilbert
representations (in particular Skinner's theorem for the local
compatibility at \(p\); see Bergdall--Hansen, Theorem 6.5.1 and
Remark 6.5.2). We take the dual of their geometric-Frobenius
representation, so arithmetic Frobenius has the raw Hecke roots
and the cyclotomic character has exponent \(+1\). The Hodge
exponents are \(b_v,n_v+b_v+1\). The eigenline of crystalline
Frobenius of slope \(-(n_v+b_v+1)\) is weakly admissible:
its filtered degree is at least \(-(n_v+b_v+1)\), and the
weak-admissibility inequality forces equality. It thus corresponds
to the high-exponent Galois line. The quotient is
\(\epsilon^{b_v}\) times the unramified character with arithmetic
Frobenius value \(p^{-b_v}\alpha_v\).
The determinant is \(\epsilon^{m+1}\xi\), first at unramified
primes and then everywhere by Chebotarev. The normalized
ordinary eigenvalues and this determinant give exactly the asserted
two residual characters on the local decomposition groups.
The residual global characteristic polynomial follows from the
Hecke congruence, the determinant, and Chebotarev and
Brauer--Nesbitt.

At \(\mathfrak q\nmid p\), the conductor exponent of the
Weil--Deligne representation is at most \(a_{\mathfrak q}(\xi)\),
by the tame \(K_0(\mathfrak n)\)-type and local newform theory.
It is at least \(a_{\mathfrak q}(\xi)\), by the determinant.
Here is why equality suffices for the exact inertia assertion.
In the Artin conductor formula the codimension of the invariants
of the determinant on every ramification subgroup is at most
that of the two-dimensional representation, and the monodromy
term is nonnegative. If the determinant is ramified, equality
forces a one-dimensional inertia-fixed space, and the two
inertial characters are \(1,\xi\); the nilpotent monodromy
commutes with inertia and is zero. If the determinant is
unramified, equality means conductor zero. Local-global
compatibility therefore gives the actual asserted inertia
representation, not just its semisimplification. At a ramified
prime the Weil representation splits into the two distinct
inertial eigenspaces, which are stable under Frobenius since
\(\xi\) is a global character. The corresponding principal series
has unlinked characters. At the remaining primes the local
representation is spherical.

Finally the cuspidal Galois representation is irreducible over
\(\overline{\mathbb Q}_p\). If its restriction to an open subgroup
is reducible, pass to the open normal core. The orbit of an
invariant line either consists of exactly two lines, permuted
transitively, or contains at least three. In the second case the
core acts by scalars, impossible because the two Hodge exponents
remain distinct after finite extension. In the first case there
is a nontrivial quadratic self-twist \(\eta\). Its reduction is
nontrivial since \(p\ne2\). Reducing the self-twist identity gives
\(\{1,\bar\xi\}=\{\bar\eta,\bar\eta\bar\xi\}\), forcing
\(\bar\xi=\bar\eta\), contrary to \(\bar\xi^2\ne1\).
\end{proof}

\paragraph{Starting eigenpacket.}
For the hypotheses above, the usual mixed weight-one Eisenstein
series supplies the starting eigenpacket. We record the argument
and its exact normalization, so that no divisibility assertion
about a special \(L\)-value is needed. Write
\(\mathfrak p_P=\prod_{v\in P}v\) and
\(\mathfrak p_Q=\prod_{v\mid p,\ v\notin P}v\).
Use the trivial ray character with modulus \(\mathfrak p_Q\)
and the ray character \(\xi\) with modulus
\(\mathfrak n\mathfrak p_P\), extending by zero at primes of
their respective moduli. Both moduli are nontrivial and the
parity condition is weight one. The Eisenstein coefficient
formula of Dasgupta--Darmon--Pollack, Proposition 2.1, gives
a form of character \(\xi\) and level
\(\mathfrak n\prod_{v\mid p}v\), with integral normalized
positive coefficients
\[
 c(\mathfrak a)=\sum_{\mathfrak d\mid\mathfrak a}
       \eta(\mathfrak a/\mathfrak d)\psi(\mathfrak d),
 \qquad c(1)=1 ,
\]
and zero constants at the standard cusps in their component
description. Its Hecke values away from \(p\mathfrak n\) and
its \(U_v\) values are precisely the desired ones, by the
coefficient formulas. We do not assert that it vanishes at
every cusp.

The weight-raising multiplier needed here is a characteristic-zero
Hilbert form of level one and trivial finite nebentypus, of positive
parallel even weight \(k\) divisible by \(p-1\), whose normalized
expansion at every standard component is \(1\bmod\lambda\).
The usual Hasse-invariant power lifting theorem supplies it:
see Goren, \emph{Lectures on Hilbert Modular Varieties and
Modular Forms}, Chapter 5, Section 6.1, Remark 6.1. For split
\(p\), use the product of the partial Hasse invariants, of
parallel weight \(p-1\), with expansion one. Equivalently
one lifts a sufficiently high power on the minimal compactification
by ampleness of the parallel Hodge bundle and Serre vanishing.
If one starts on polarized components, the power can be taken
even and the lift descended with trivial character by taking
the product of its translates under the finite component and
level descent group, acting with the weight-normalized central
action. The Hasse section is invariant for that action; hence
the product still has expansion one modulo the uniformizer.
This norm construction does not divide by the order of a class
group. In the normalized ideal-coefficient convention one chooses
the component ideals prime to \(p\) and takes \(k\) divisible by
\(2(p-1)\), so the factors \(N(\mathfrak t)^{\pm k/2}\)
are also one modulo \(p\).

Multiplication by this lift gives a nonzero residual eigenpacket
at weight \(k+1\), at the same level and character, since the
positive coefficient at \(1\) remains one. The Hecke coefficient
formulas are compatible modulo \(\lambda\): at primes away from
\(p\mathfrak n\) the extra factor is \(N\mathfrak q^k\equiv1\),
and at \(p\) the \(U_v\) formula is the coefficient shift.
Use the lattice of forms with integral normalized coefficients
at the standard components. It is finite over \(O\) by the
\(q\)-expansion principle (finitely many coefficients separate
the finite-dimensional space) and stable under these operators.
The Deligne--Serre lemma lifts the residual eigenpacket in the
full holomorphic space.

The lift cannot be Eisenstein. In weight \(k+1>2\) an Eisenstein
eigensystem at this character and level is obtained from two
finite characters \(\eta',\psi'\) with product \(\xi\), and its
unramified roots are \(\eta'(\mathfrak q)\) and
\(\psi'(\mathfrak q)N\mathfrak q^k\). The residual unordered
characters are \(\{1,\bar\xi\}\) by Chebotarev and
Brauer--Nesbitt. At \(v\mid p\) both characters are unramified:
their conductor exponents have sum at most one, and their
product is unramified. The only unit Hecke root there is
\(\eta'(v)\). It cannot reduce to \(1\) on \(P\) and to
\(\bar\xi(v)\) on its complement, as both sets are nonempty and
\(\bar\xi(v)\ne1\) at every such prime. Thus the cuspidal part
contains the lifted system. This uses separation of ordinary
refinements, not cuspidality of the weight-one series itself.

\paragraph{Source normalization and scope.}
For an arithmetic Hecke polynomial with roots \(\alpha,\beta\)
the normalized Satake parameters of \(\pi_0\) are
\(q^{-1/2}(\alpha,\beta)\). The source's attachment uses inverse
(arithmetic) Frobenius and multiplies its normalized Satake
parameters by \(q^{1/2}\). Accordingly take
\(\pi=\pi_0^\vee\otimes|\det|\), with normalized parameters
\(q^{-1/2}(\alpha^{-1},\beta^{-1})\).
At a real place the quaternion coefficient representation has
highest weights \(n_v+b_v,b_v\). The Harish--Chandra parameters
of \(\pi_0\) are
\(\{-b_v+1/2,-n_v-b_v-1/2\}\); dualizing and twisting as above
gives \(\{b_v+1/2,n_v+b_v+3/2\}\).
These are exactly the source's rank-two \(\operatorname{infOf}\)
parameters for the two cyclotomic exponents. Principal series
with distinct inertia characters are unlinked, and the other
places have a Satake parameter. Thus the construction targets
the original \texttt{MixSeedP} interface, with its original
field and local orders, not a parallel-weight substitute.

The integral weight-change theorem is independent of the
Hasse multiplier. The only additional input in the starting
packet argument is the classical level-one Hasse-power lift
in the normalized adelic convention described above and the
weight-one Eisenstein coefficient formula. In particular this
does not prove the subsequent symplectic induction, the second
or third case of rank-four propagation, or the final uniform
Selmer estimate.

\paragraph{References.}
Dasgupta--Darmon--Pollack, \emph{Hilbert modular forms and the
Gross--Stark conjecture}, Annals of Mathematics 174 (2011),
439--484, Proposition 2.1 (including imprimitive moduli);
Goren, \emph{Lectures on Hilbert Modular Varieties and Modular
Forms}, CRM Monograph Series 14, 2002, Chapter 5,
Sections 3 and 6.1;
Skinner, \emph{A note on the \(p\)-adic Galois representations
attached to Hilbert modular forms}, Documenta Mathematica 14
(2009), 241--258;
Bergdall--Hansen, \emph{On \(p\)-adic \(L\)-functions for Hilbert
modular forms}, Theorem 6.5.1 and Remark 6.5.2.
